\subsection{Discussion}
\label{sec:discussion}

The results confirm the central hypothesis: a protocol-aware DRL agent trained with domain
randomisation and curriculum learning outperforms standard path-planning heuristics across a
range of network densities, with the advantage growing as the problem becomes harder.

\subsubsection{Mechanism of the DQN Advantage}

The DQN's key advantage is its ability to learn implicit positioning rules that maximise
concurrent transmission across orthogonal spreading factors. Per-step action analysis confirms
that greedy agents issue a collection action at every timestep (100\% hover fraction), whereas
the DQN devotes 2.9\% of steps to deliberate movement ($p = 0.002$, Welch's $t$-test, $n=10$
seeds, Cohen's $d = 2.1$, large effect), reaching positions yielding $48.7$\,B per collection
step versus $42.7$\,B for SF-Aware Greedy (+14\%). Capture Effect trigger counts are
statistically indistinguishable ($p = 1.0$, $d \approx 0$), confirming the gain is driven by
improved spatial positioning, not by resolving intra-SF collisions.

The SF dynamics analysis (Figure~\ref{fig:sf_dynamics}) provides the causal explanation:
by maintaining proximity for multiple steps, the DQN allows the EMA-ADR filter to converge
to lower SF assignments before repositioning. This converts ADR latency---normally a fidelity
cost---into a feature: induced SF downgrades enable high-throughput orthogonal-SF concurrent
collection that greedy policies cannot plan for.

\subsubsection{Density Dependence and Effect Sizes}

All four density conditions reach statistical significance.
At $N=10$, $d = 2.52$ (large effect, $p < 0.001$): both agents have very tight reward
distributions at low density, so the mean gap ($+2.3\%$) is highly reliably reproduced.
At $N=20$, $d = 0.66$ (medium effect, $p = 0.049$) and at $N=30$, $d = 0.78$
(medium effect, $p = 0.021$). At $N=40$, extending to thirty seeds yields $+8.6\%$
($d = 0.54$, medium effect, $p = 0.043$), confirming significance at the highest density.
The non-monotonic pattern ($d=0.54$ at $N=40$ vs.\ $d=0.78$ at $N=30$) reflects
layout-dependent variance dominating the mean gap at extreme density (std = 1.120M at
$N=40$ vs.\ 0.040M at $N=10$). Table~\ref{tab:effect_sizes} summarises all effect sizes.

\begin{table}[H]
\centering
\caption{Cohen's $d$ effect sizes for DQN vs.\ SF-Aware Greedy (cumulative reward),
         20-seed evaluation (30-seed for $N=40$). $d < 0.2$ = negligible; $0.2$--$0.5$ = small;
         $0.5$--$0.8$ = medium; $> 0.8$ = large.}
\label{tab:effect_sizes}
\begin{tabularx}{\textwidth}{l Z Z Z Z}
\toprule
$N$ & Reward $d$ & Jain's $d$ & Efficiency $d$ & Interpretation \\
\midrule
10  & 2.52 & 0.43 & 2.68 & Large ($p < 0.001$) \\
20  & 0.66 & 0.36 & 0.67 & Medium ($p = 0.049$\textsuperscript{*})  \\
30  & 0.78 & 0.61 & 0.77 & Medium ($p = 0.021$\textsuperscript{*})  \\
40  & 0.54 & 0.40 & 0.55 & Medium ($p = 0.043$\textsuperscript{*}) \\
\bottomrule
\end{tabularx}
\end{table}

\subsubsection{Limitations}

Two limitations apply. First, fairness (Jain's Index) at $N=40$ does not reach significance ($p=0.129$), degrading for all agents due to the fixed battery constraint, not the algorithm~\cite{jin2023equalizing}. Second, all evaluation is simulation-based; real-world effects (wind, GPS error, hardware latency) are discussed in Chapter~\ref{sec:conclusions}.
